% \vspace{-0.2in}
\section{Related Works}

% \vspace{-0.12in}
\paragraph{From Multi-Step Retrieval to Interactive Search Agents}
Prior work in information retrieval has largely centered on single-shot retrieve-and-read pipelines \citep{lee-etal-2019-latent, guu2020retrieval, karpukhin-etal-2020-dense}. This paradigm expanded with the emergence of LLM-based systems that perform retrieval over multiple steps. Prompted tool-use frameworks such as ReAct \citep{react} and multi-step retrieval methods such as IRCoT \citep{ircot} showed that LLMs can improve retrieval coverage on complex queries by iteratively issuing retrieval actions conditioned on intermediate reasoning or generations. These approaches established the broader view that retrieval can be formulated as a sequential, interactive process rather than a single retrieval step.
Building on this perspective, subsequent work increasingly developed interactive search agents that actively generate queries, consume observations, and iteratively refine their information needs over the course of search. Search-o1 \citep{search-o1} demonstrated the effectiveness of this search-centric agent design through prompting-based iterative query formulation and evidence collection over the web. More recent approaches, including \citet{search-r1}, \citet{r1-searcher}, \citet{deepresearcher}, and \citet{ReSearch}, have further advanced this paradigm by training search agents with reinforcement learning \citep{ppo, grpo}. In these approaches, search is explicitly modeled as sequential decision-making over query actions and observation consumption, making RL a natural framework for optimizing end-to-end search behavior.

% \vspace{-0.12in}
\paragraph{Beyond Gold Supervision for Search Agents}
Many RL-based search agents have demonstrated strong performance, but they remain dependent on gold supervision, such as ground-truth answers, to construct reliable reward signals. This dependence creates a scalability bottleneck in domains where high-quality annotations are scarce, expensive, or rapidly evolving.
To reduce the cost of annotation, several recent methods have explored alternatives to gold supervision by leveraging alternative supervision signals, including intrinsic signals and self-generated pseudo-labels. RLIF \citep{rlif}, for example, derives reward signals from the model's internal confidence to guide training without external labels.
Rubric-based LLM judges \citep{geval}, which we refer to as Constitutional Judges, can also be used to provide reward signals by scoring model outputs according to a predefined set of evaluation criteria.
Test-time RL methods such as TTRL \citep{ttrl} instead use agreement among multiple sampled rollouts as a proxy signal for updating policies on unlabeled instances. While these approaches reduce reliance on human annotation, their training signals are not specifically aligned with the central objective of search: rewarding trajectories that acquire sufficient and necessary external evidence to satisfy the original information need.
% Co-evolutionary or self-play style frameworks (e.g., Dr. Zero \citep{drzero}) further explore data-free improvement loops with automated curricula.

% \vspace{-0.12in}
\paragraph{Extending Cycle-Consistency to Search Trajectories}
Cycle-consistency has been widely used as a learning principle that encourages transformations to preserve essential information, notably in unsupervised machine translation \citep{he2016dual, lample2017unsupervised, han2021unsupervised} and image-to-image translation \citep{cyclegan, liu2017unsupervised, huang2018multimodal}. Its central intuition is that a good transformation should retain enough information for the original input to be recoverable through an inverse mapping. Related reconstruction-based objectives have also been explored in retrieval-augmented pretraining, where retrieved documents serve as latent variables for reconstructing targets \citep{guu2020retrieval}. In the search setting, we consider cycle-consistency over agentic search trajectories, treating a trajectory as an information-preserving encoding of the source question. A central challenge in this setting is information leakage: lexical overlap between the question and intermediate search queries (or the final response) can enable trivial reconstruction without reflecting genuine search quality. We address this issue through information bottlenecks, so that successful reconstruction must rely on retrieved observations together with the structural scaffold rather than surface-form redundancy.